\textbf{Phần 1. Trắc nghiệm lựa chọn}
\setcounterpageref{ex}{0}

\begin{ex} %[5D1N6]
Kết quả của phép cộng $\dfrac{1}{7} + \dfrac{1}{9}$ là
\choice
{$\dfrac{2}{16}$}
{$\dfrac{2}{63}$}
{\True $\dfrac{16}{63}$}
{$\dfrac{14}{63}$}
\loigiai{
Quy đồng mẫu số chung là $63$: $\dfrac{1}{7} + \dfrac{1}{9} = \dfrac{9}{63} + \dfrac{7}{63} = \dfrac{16}{63}$.\\
\textbf{Đáp án của câu hỏi là C}
}
\end{ex}

\begin{ex} %[5D1N6]
Kết quả của phép trừ $\dfrac{1}{4} - \dfrac{1}{9}$ là
\choice
{\True $\dfrac{5}{36}$}
{$\dfrac{1}{5}$}
{$\dfrac{13}{36}$}
{$\dfrac{5}{13}$}
\loigiai{
Quy đồng mẫu số chung $36$: $\dfrac{1}{4} - \dfrac{1}{9} = \dfrac{9}{36} - \dfrac{4}{36} = \dfrac{5}{36}$.\\
\textbf{Đáp án của câu hỏi là A}
}
\end{ex}

\begin{ex} %[5D1H6]
Kết quả của phép cộng $\dfrac{3}{11} + \dfrac{7}{12}$ là
\choice
{$\dfrac{10}{23}$}
{\True $\dfrac{113}{132}$}
{$\dfrac{10}{132}$}
{$\dfrac{111}{132}$}
\loigiai{
Quy đồng mẫu số chung là $132$:
\[
\dfrac{3}{11} + \dfrac{7}{12} = \dfrac{3 \times 12}{132} + \dfrac{7 \times 11}{132} = \dfrac{36 + 77}{132} = \dfrac{113}{132}.
\]
\textbf{Đáp án của câu hỏi là B}
}
\end{ex}

\begin{ex} %[5D1H6]
Kết quả của phép trừ $\dfrac{2}{5} - \dfrac{1}{4}$ là
\choice
{$\dfrac{1}{1}$}
{$\dfrac{1}{20}$}
{$\dfrac{7}{20}$}
{\True $\dfrac{3}{20}$}
\loigiai{
Quy đồng mẫu số chung là $20$: $\dfrac{2}{5} - \dfrac{1}{4} = \dfrac{8}{20} - \dfrac{5}{20} = \dfrac{3}{20}$.\\
\textbf{Đáp án của câu hỏi là D}
}
\end{ex}

\begin{ex} %[5D1H6]
Kết quả của phép cộng $\dfrac{6}{11} + \dfrac{1}{2}$ là
\choice
{\True $\dfrac{23}{22}$}
{$\dfrac{7}{13}$}
{$\dfrac{13}{22}$}
{$\dfrac{17}{22}$}
\loigiai{
Quy đồng mẫu số: $\dfrac{6}{11} + \dfrac{1}{2} = \dfrac{12}{22} + \dfrac{11}{22} = \dfrac{23}{22}$.\\
\textbf{Đáp án của câu hỏi là A}
}
\end{ex}

\begin{ex} %[5D1H6]
Rút gọn rồi tính hiệu $\dfrac{20}{18} - \dfrac{2}{5}$ ta được kết quả là
\choice
{$\dfrac{18}{13}$}
{\True $\dfrac{32}{45}$}
{$\dfrac{14}{45}$}
{$\dfrac{16}{45}$}
\loigiai{
Rút gọn $\dfrac{20}{18} = \dfrac{10}{9}$.\\
Thực hiện phép trừ: $\dfrac{10}{9} - \dfrac{2}{5} = \dfrac{50}{45} - \dfrac{18}{45} = \dfrac{32}{45}$.\\
\textbf{Đáp án của câu hỏi là B}
}
\end{ex}

\begin{ex} %[5D1H6]
Dấu thích hợp điền vào ô trống để $\dfrac{1}{20} \dots \dfrac{2}{3} = \dfrac{43}{60}$ là
\choice
{$-$}
{$\times$}
{$:$}
{\True $+$}
\loigiai{
Ta có: $\dfrac{1}{20} + \dfrac{2}{3} = \dfrac{3}{60} + \dfrac{40}{60} = \dfrac{43}{60}$. Vậy dấu cần điền là dấu $+$.\\
\textbf{Đáp án của câu hỏi là D}
}
\end{ex}

\begin{ex} %[5D1H6]
Dấu thích hợp điền vào ô trống để $\dfrac{7}{11} \dots \dfrac{5}{8} = \dfrac{1}{88}$ là
\choice
{$+$}
{\True $-$}
{$\times$}
{$:$}
\loigiai{
Ta có: $\dfrac{7}{11} - \dfrac{5}{8} = \dfrac{56}{88} - \dfrac{55}{88} = \dfrac{1}{88}$. Vậy dấu cần điền là dấu $-$.\\
\textbf{Đáp án của câu hỏi là B}
}
\end{ex}

\begin{ex} %[5D1H6]
Kết quả của phép tính $7 + \dfrac{1}{5}$ dưới dạng phân số là
\choice
{$\dfrac{8}{5}$}
{$\dfrac{7}{5}$}
{\True $\dfrac{36}{5}$}
{$\dfrac{35}{5}$}
\loigiai{
Viết $7 = \dfrac{35}{5}$. Khi đó $7 + \dfrac{1}{5} = \dfrac{35}{5} + \dfrac{1}{5} = \dfrac{36}{5}$.\\
\textbf{Đáp án của câu hỏi là C}
}
\end{ex}

\begin{ex} %[5D1H6]
Kết quả của phép tính $3 - \dfrac{11}{8}$ dưới dạng phân số là
\choice
{$\dfrac{8}{8}$}
{$\dfrac{14}{8}$}
{$\dfrac{21}{8}$}
{\True $\dfrac{13}{8}$}
\loigiai{
Viết $3 = \dfrac{24}{8}$. Khi đó $3 - \dfrac{11}{8} = \dfrac{24}{8} - \dfrac{11}{8} = \dfrac{13}{8}$.\\
\textbf{Đáp án của câu hỏi là D}
}
\end{ex}

\begin{ex} %[5D1V6]
Có hai vòi nước cùng chảy vào một bể cạn. Vòi thứ nhất chảy được $\dfrac{1}{2}$ bể, vòi thứ hai chảy được $\dfrac{2}{5}$ bể. Cả hai vòi đã chảy được bao nhiêu phần bể nước?
\choice
{$\dfrac{3}{7}\text{ bể}$}
{$\dfrac{3}{10}\text{ bể}$}
{\True $\dfrac{9}{10}\text{ bể}$}
{$1\text{ bể}$}
\loigiai{
Cả hai vòi chảy được số phần bể nước là: $\dfrac{1}{2} + \dfrac{2}{5} = \dfrac{5}{10} + \dfrac{4}{10} = \dfrac{9}{10}$ (bể).\\
\textbf{Đáp án của câu hỏi là C}
}
\end{ex}

\begin{ex} %[5D1V6]
Để chia đều $5$ chiếc bánh cho $6$ người, mỗi người sẽ nhận được bao nhiêu phần chiếc bánh?
\choice
{\True $\dfrac{5}{6}\text{ chiếc bánh}$}
{$\dfrac{6}{5}\text{ chiếc bánh}$}
{$\dfrac{1}{6}\text{ chiếc bánh}$}
{$\dfrac{1}{5}\text{ chiếc bánh}$}
\loigiai{
Chia đều $5$ chiếc bánh cho $6$ người: mỗi người được $5 : 6 = \dfrac{5}{6}$ chiếc bánh.\\
\textbf{Đáp án của câu hỏi là A}
}
\end{ex}

\textbf{Phần 2. Trắc nghiệm đúng sai}
\setcounterpageref{ex}{0}

\begin{ex} %[5D1H6]
Cho các phép tính: $P = \dfrac{13}{5} + \dfrac{7}{6}$, $Q = \dfrac{11}{7} - \dfrac{5}{8}$ và $R = \dfrac{11}{4} - \dfrac{7}{5}$. Xét tính đúng sai của các khẳng định sau:
\choiceTF
{\True Phép tính $P$ có kết quả là $\dfrac{113}{30}$.}
{Phép tính $Q$ có kết quả là $\dfrac{6}{56}$.}
{\True Phép tính $R$ có kết quả là $\dfrac{27}{20}$.}
{\True Kết quả phép tính $P$ lớn hơn $3$.}
\loigiai{
\begin{itemize}
    \item $P = \dfrac{13 \times 6 + 7 \times 5}{30} = \dfrac{78 + 35}{30} = \dfrac{113}{30}$.
    \item $Q = \dfrac{11 \times 8 - 5 \times 7}{56} = \dfrac{88 - 35}{56} = \dfrac{53}{56} \neq \dfrac{6}{56}$.
    \item $R = \dfrac{11 \times 5 - 7 \times 4}{20} = \dfrac{55 - 28}{20} = \dfrac{27}{20}$.
    \item Vì $\dfrac{113}{30} = 3\dfrac{23}{30} > 3$.
\end{itemize}
\begin{itemize}
    \item a) \textbf{Đúng}.
    \item b) \textbf{Sai}: $Q = \dfrac{53}{56}$.
    \item c) \textbf{Đúng}.
    \item d) \textbf{Đúng}.
\end{itemize}
}
\end{ex}

\begin{ex} %[5D1V6]
Việt đổ $\dfrac{1}{5}\text{ l}$ nước vào bình, Mai đổ $\dfrac{1}{2}\text{ l}$ nước vào cùng chiếc bình đó. Xét tính đúng sai của các khẳng định sau:
\choiceTF
{\True Cả hai bạn đã đổ tất cả $\dfrac{7}{10}\text{ l}$ nước vào bình.}
{Mai đổ ít nước hơn Việt.}
{\True Mai đổ nhiều hơn Việt $\dfrac{3}{10}\text{ l}$ nước.}
{Nếu dung tích của bình là $1\text{ l}$ thì lượng nước còn thiếu để đầy bình là $\dfrac{1}{3}\text{ l}$.}
\loigiai{
\begin{itemize}
    \item Cả hai bạn đổ: $\dfrac{1}{5} + \dfrac{1}{2} = \dfrac{2}{10} + \dfrac{5}{10} = \dfrac{7}{10}\text{ (l)}$.
    \item Vì $\dfrac{1}{2} = \dfrac{5}{10} > \dfrac{2}{10} = \dfrac{1}{5}$ nên Mai đổ nhiều hơn Việt.
    \item Lượng nước Mai đổ nhiều hơn Việt là: $\dfrac{1}{2} - \dfrac{1}{5} = \dfrac{3}{10}\text{ (l)}$.
    \item Lượng nước còn thiếu để đầy bình $1\text{ l}$ là: $1 - \dfrac{7}{10} = \dfrac{3}{10}\text{ (l)} \neq \dfrac{1}{3}\text{ l}$.
\end{itemize}
\begin{itemize}
    \item a) \textbf{Đúng}.
    \item b) \textbf{Sai}: Mai đổ nhiều hơn Việt.
    \item c) \textbf{Đúng}.
    \item d) \textbf{Sai}: Còn thiếu $\dfrac{3}{10}\text{ l}$.
\end{itemize}
}
\end{ex}

\textbf{Phần 3. Trả lời ngắn}
\setcounterpageref{ex}{0}

\begin{ex} %[5D1H6]
Tính tổng $\dfrac{3}{10} + \dfrac{5}{21}$ dưới dạng phân số tối giản $\dfrac{a}{b}$. Giá trị của tử số $a$ bằng bao nhiêu?
\shortans{113}
\loigiai{
\textbf{Đáp số:} $113$\\
$\dfrac{3}{10} + \dfrac{5}{21} = \dfrac{3 \times 21 + 5 \times 10}{210} = \dfrac{63 + 50}{210} = \dfrac{113}{210}$. Tử số $a = 113$.
}
\end{ex}

\begin{ex} %[5D1H6]
Tính hiệu $\dfrac{31}{10} - 2$ dưới dạng phân số tối giản $\dfrac{m}{n}$. Tính giá trị của $m + n$.
\shortans{21}
\loigiai{
\textbf{Đáp số:} $21$\\
$\dfrac{31}{10} - 2 = \dfrac{31}{10} - \dfrac{20}{10} = \dfrac{11}{10}$. Khi đó $m = 11, n = 10 \Rightarrow m + n = 21$.
}
\end{ex}

\begin{ex} %[5D1H6]
Tính tổng $\dfrac{14}{5} + 1$ dưới dạng phân số tối giản $\dfrac{p}{q}$. Tìm giá trị của tử số $p$.
\shortans{19}
\loigiai{
\textbf{Đáp số:} $19$\\
$\dfrac{14}{5} + 1 = \dfrac{14}{5} + \dfrac{5}{5} = \dfrac{19}{5}$. Tử số $p = 19$.
}
\end{ex}

\begin{ex} %[5D1V6]
Có hai vòi nước cùng chảy vào một bể cạn. Vòi thứ nhất chảy được $\dfrac{1}{2}$ bể, vòi thứ hai chảy được $\dfrac{2}{5}$ bể. Hỏi còn lại bao nhiêu phần bể chưa có nước? (Viết kết quả dưới dạng phân số tối giản $\dfrac{1}{k}$). Giá trị của $k$ là bao nhiêu?
\shortans{10}
\loigiai{
\textbf{Đáp số:} $10$\\
Cả hai vòi chảy được: $\dfrac{1}{2} + \dfrac{2}{5} = \dfrac{9}{10}$ (bể).\\
Phần bể chưa có nước là: $1 - \dfrac{9}{10} = \dfrac{1}{10}$ (bể). Vậy $k = 10$.
}
\end{ex}

\begin{ex} %[5D1V6]
Mai và Việt cùng đổ nước vào một bình. Mai đổ $\dfrac{1}{2}\text{ l}$, Việt đổ $\dfrac{1}{5}\text{ l}$. Lượng nước Mai đổ nhiều hơn Việt là $\dfrac{x}{10}\text{ l}$. Giá trị của $x$ là bao nhiêu?
\shortans{3}
\loigiai{
\textbf{Đáp số:} $3$\\
$\dfrac{1}{2} - \dfrac{1}{5} = \dfrac{5}{10} - \dfrac{2}{10} = \dfrac{3}{10}\text{ (l)} \Rightarrow x = 3$.
}
\end{ex}

\begin{ex} %[5D1V6]
Rô-bốt chia bánh cho bạn bằng cách lấy $\dfrac{1}{2}$ cái bánh cộng thêm $\dfrac{1}{3}$ cái bánh. Tổng số bánh bạn nhận được là $\dfrac{u}{6}$ cái bánh. Giá trị của $u$ là bao nhiêu?
\shortans{5}
\loigiai{
\textbf{Đáp số:} $5$\\
$\dfrac{1}{2} + \dfrac{1}{3} = \dfrac{3}{6} + \dfrac{2}{6} = \dfrac{5}{6}\text{ (cái bánh)} \Rightarrow u = 5$.
}
\end{ex}
